\clearpage
\section{051 -- Zone escluse dalle regole clearance ECSS generiche}

\begin{tabularx}{\linewidth}{@{}>{\bfseries}l X@{}}
\toprule
Data & 13/11/2026 (problema segnalato il 12/11/2026)\\
Revisione & non indicata\\
Redatto da & Claude Code\\
Approvazione & Brik81\\
Stato & \tg{tagdone}{MODIFICA ESEGUITA} \quad \tg{tagwarn}{CRITICITÀ APERTE}\\
\bottomrule
\end{tabularx}

Solo regole DRC: nessun componente, nessuna net, nessun codice JLCPCB/Mouser toccati. Nessun altro file modificato oltre a \code{SYMPHONY.kicad\_dru}.

\subsection{Componenti, net e fogli coinvolti}
\begin{itemize}
  \item File modificato: \code{SYMPHONY.kicad\_dru}.
  \item Regole: ``ECSS esterni clearance 0.20'' (layer esterni) e ``ECSS interni clearance 0.13'' (layer interni).
  \item Zone d'esempio su \code{In5.Cu}: \code{3V3} (clearance zona \SI{1,0}{\milli\metre}, priorità 1), \code{5V0} (\SI{0,5}{\milli\metre}, priorità 3), \code{5V} (\SI{0,2}{\milli\metre}, priorità 2).
  \item Zone RF su \code{F.Cu} (clearance zona 0, quindi segue la net class): \code{PA\_VCC2}, \code{RF\_PA\_IN}, \code{RF\_PA\_OUT}.
\end{itemize}

\subsection{Problema segnalato}
Brik81 ha notato che, cambiando la clearance delle zone in KiCad e rifacendo il riempimento, le distanze tra zone di net diverse non cambiavano e restavano bassissime.

Analisi: il riempimento è stato rifatto con pcbnew di KiCad 10 su una copia del progetto e sono state misurate le distanze tra i riempimenti su \code{In5.Cu}.

\begin{center}
\small
\begin{tabular}{@{}lccc@{}}
\toprule
Coppia di zone (\code{In5.Cu}) & Clearance impostata & Misurata prima & Misurata dopo\\
\midrule
\code{3V3}--\code{5V0} & \SI{1,0}{\milli\metre} & \SI{0,135}{\milli\metre} & \SI{1,002}{\milli\metre}\\
\code{3V3}--\code{5V} & \SI{1,0}{\milli\metre} & \SI{0,134}{\milli\metre} & \SI{1,002}{\milli\metre}\\
\code{5V0}--\code{5V} & \SI{0,5}{\milli\metre} & \SI{0,135}{\milli\metre} & \SI{0,501}{\milli\metre}\\
\bottomrule
\end{tabular}
\end{center}

Causa: le due regole ECSS generiche non avevano condizioni. In KiCad una regola custom prevale sulla clearance propria della zona e sulla clearance della net class (non si prende il valore massimo). Quindi su tutti i layer interni, zona--zona e zona--oggetto, valevano \SI{0,13}{\milli\metre}, e sugli esterni \SI{0,20}{\milli\metre}, qualunque fosse il valore impostato nella zona.

\subsection{Cosa è cambiato}
\begin{center}
\small
\begin{tabularx}{\linewidth}{@{}l l X X@{}}
\toprule
Regola & Layer & Prima & Dopo\\
\midrule
ECSS esterni clearance 0.20 & outer & nessuna condizione; min \SI{0,2}{\milli\metre} & \code{(condition "A.Type != 'Zone' \&\& B.Type != 'Zone'")}; min \SI{0,2}{\milli\metre} invariato\\
ECSS interni clearance 0.13 & inner & nessuna condizione; min \SI{0,13}{\milli\metre} & stessa condizione; min \SI{0,13}{\milli\metre} invariato\\
\bottomrule
\end{tabularx}
\end{center}

\begin{itemize}
  \item Aggiunto in \code{SYMPHONY.kicad\_dru} un commento datato 13/11/2026 che spiega il motivo.
  \item Piste, via e pad tra loro restano sotto le regole ECSS come prima.
  \item Per le zone vale ora la clearance della zona, oppure quella della net class se la zona ha clearance 0.
  \item Invariate le altre regole con condizione: RF \SI{0,35}{\milli\metre}, HV\_50V, via-via 0.13/0.20, \code{hole\_clearance}, \code{edge\_clearance}.
\end{itemize}

\textbf{Verifica fatta} (riempimento rifatto con pcbnew su una copia, con il nuovo \code{.dru}):
\begin{itemize}
  \item \code{In5.Cu}: valori nella tabella sopra; prima erano tutte \SI{0,13}{\milli\metre}, ora le clearance impostate sono rispettate.
  \item \code{F.Cu}, zone RF che seguono la net class: \code{PA\_VCC2}--\code{RF\_PA\_IN} \SI{1,088}{\milli\metre}, \code{PA\_VCC2}--\code{RF\_PA\_OUT} \SI{1,220}{\milli\metre}, \code{RF\_PA\_IN}--\code{RF\_PA\_OUT} \SI{2,297}{\milli\metre}.
\end{itemize}

\subsection{Motivazione}
Sono state proposte due opzioni a Brik81:
\begin{enumerate}
  \item Escludere solo le zone dalle due regole ECSS generiche (consigliata, cambiamento minimo).
  \item Limitare le regole generiche alla net class \code{Default}, così da far valere anche le clearance delle net class; più invasiva, con possibili nuove violazioni DRC sulle piste RF.
\end{enumerate}
Brik81 ha scelto la prima.

Approvazione di Brik81 (citazione testuale):
\begin{quote}
``si escludi le zone dalle regole''
\end{quote}

\subsection{Impatti e criticità}
\begin{critbox}
\textbf{Riempimento da rifare.} Il \code{.kicad\_pcb} sul disco contiene ancora i riempimenti vecchi a \SI{0,13}{\milli\metre}. Brik81 deve rifare il riempimento di tutte le zone in KiCad (tasto B) e salvare.
\end{critbox}
\begin{critbox}
\textbf{Frammentazione dei piani.} Le zone con clearance grande (\code{3V3} a \SI{1,0}{\milli\metre}, GND interni a \SI{0,8}{\milli\metre}) lasciano ora attorno a via e pad di altre net vuoti più larghi. Il piano GND su In1/In3/In6 potrebbe risultare troppo frammentato, con nuove isole o colli stretti (\code{min\_thickness}).
\end{critbox}
\begin{notebox}
\textbf{Minimo ECSS sulle zone esterne.} Su F.Cu/B.Cu viene meno il minimo di \SI{0,20}{\milli\metre} ECSS per le zone. Con clearance zona 0 si usa la net class (\code{Default} \SI{0,20}{\milli\metre}, quindi equivalente). Una zona con clearance locale impostata sotto \SI{0,20}{\milli\metre} su F.Cu/B.Cu non sarebbe più segnalata dalla regola ECSS.
\end{notebox}
\begin{critbox}
\textbf{Aperto per scelta di Brik81.} Le clearance più grandi delle net class \code{RF\_CLK} (\SI{0,30}{\milli\metre}), \code{RF\_PA\_OUT} (\SI{0,25}{\milli\metre}) e \code{RF\_50R\_CPWG} (\SI{0,25}{\milli\metre}) non hanno effetto su piste e via, che restano a \SI{0,20}{\milli\metre} sugli esterni e \SI{0,13}{\milli\metre} sugli interni. Brik81 ha scelto di non intervenire per ora.
\end{critbox}

\subsection{Verifiche da fare}
\begin{itemize}
  \item Rifare il riempimento di tutte le zone (tasto B) e salvare il progetto.
  \item Lanciare il DRC dopo il riempimento.
  \item Controllare la frammentazione del GND su In1/In3/In6 e l'assenza di nuove isole o colli stretti (\code{min\_thickness}).
  \item Tenere presente che zone esterne con clearance locale sotto \SI{0,20}{\milli\metre} non sono più segnalate.
  \item Decidere in futuro se far valere le clearance delle net class RF su piste e via (opzione 2).
\end{itemize}
